\subsection{Drain (DRN) Package}

The Drain Package removes water from the aquifer at a rate proportional to the
amount by which the aquifer head exceeds a specified drain elevation, and removes
no water when the head is below that elevation. Its two per-boundary inputs are
the drain elevation and the conductance.

\begin{longtable}{p{0.15\textwidth} p{0.15\textwidth} p{0.60\textwidth}}
\caption{Memory-manager variables in the DRN Package that are candidates for
modification through the API. See also the variables common to all stress
packages described above.} \label{table:api-gwf-drn} \\

\hline
\textbf{Variable} & \textbf{Class} & \textbf{Guidance} \\
\hline
\endfirsthead

\hline
\textbf{Variable} & \textbf{Class} & \textbf{Guidance} \\
\hline
\endhead

\texttt{ELEV} & Period-reset &
Drain elevation for each drain, dimensioned by \texttt{MAXBOUND}. Read from input
at the start of each stress period; set it every time step to change the drain
elevation at run time. \\

\texttt{COND} & Period-reset &
Drain conductance for each drain, dimensioned by \texttt{MAXBOUND}. Read from
input at the start of each stress period; set it every time step to change the
conductance at run time. When a drain-depth auxiliary variable
(\texttt{AUXDEPTHNAME}) is used, the conductance applied as the head approaches
the drain elevation is scaled from \texttt{COND} (linearly under the standard
formulation, cubically under the Newton-Raphson formulation); \texttt{COND}
remains the unscaled conductance to set. \\

\texttt{AUXVAR} & Period-reset &
Auxiliary variable values, dimensioned by the number of auxiliary variables and
\texttt{MAXBOUND} (present when \texttt{AUXILIARY} variables are specified).
Re-read each stress period; set after \texttt{prepare\_time\_step}. Indexed by
boundary position, so a value follows the boundary index, not the cell. \\

\texttt{NODELIST} & Period-reset &
One-based reduced node number of the cell containing each boundary, dimensioned
by \texttt{MAXBOUND}. Set it after \texttt{prepare\_time\_step} to relocate a
boundary (see ``Changing Boundary Locations''). \\

\texttt{NBOUND} & Period-reset &
Number of active boundaries in the current stress period, from 0 to
\texttt{MAXBOUND}. Increase it (and fill the new entries) to activate boundaries
or decrease it to deactivate them (see ``Changing Boundary Locations''). \\

\hline
\end{longtable}

\noindent As with all of these variables, a value supplied by a time series is
refreshed each time step and overwrites API changes.
